\section{Hardware in context}

Tiny Tapeout is a development platform for Academics, Industry, and Hobbyests alike to develop custom silicon designs for a variety of applications.
Their system is fully open source and is based on open source PDKs and open development tools such as XScheme and MAGIC VLSI. 
The core concept behind Tiny Tapeout is an extension of earlier multi-project wafer (MPW) services such as those offered by organizations such as CMC. 
Instead of ordering a wafer with multiple projects on it like in MPW, they are instead a participant in an MPW run and instead load multiple projects onto a single die within that MPW run. 
This substantially lowers the cost bringing it into the realm of what is obtainable for an individual or low-budget labratory project. 
They package and return a piece of silicon containing many peoples projects, with associated select circuitry to allow 8 inputs, 8 outputs, and 8 bi-direction to be routed to the outside of the chip.
They also have a limited number of spaces which can also contain up to 8 analog pins; however, they are not relevant for our purposes. 
Those 24 input outputs are routed to a set of 3 PMOD interfaces. 
The board described here interacts with those 3 PMOD interfaces to allow for a combination of potential routing, and capture/ stimulation of those values. 


To test a design one may typically add external circuits to validate that everything functions as expected. 
This may involve both creaitng expected stimulus and capturing the output to validate it against simulation. 
This would then require both a waveform generator and a digital logic analyzer. 
The board here acts as an interposer. 
This way the outputs and inputs can be captured without requiring additional tools. 
They can also be captured while interacting with the expected output hardware. 